\paragraph{Table 1 of the article.} \emph{The 22 decision layers of release 1.1.0 by group, and the two layers of Daisy output.}

\textbf{How it is made.} Descriptive table (definitions, design, parameters or tests); no computation.

\paragraph{Table 2 of the article.} \emph{Questions of the workbench and the layers that answer them.}

\textbf{How it is made.} Descriptive table (definitions, design, parameters or tests); no computation.

\paragraph{Table 3 of the article.} \emph{Soils used for the test of the soil column (van Genuchten parameters, $\ell=0.5$).}

\textbf{How it is made.} Table typed from the result file (13 of its numbers found there). Script: \nolinkurl{Manuscripts/reproducibility/p3_js/run_drip.js}. Input: \nolinkurl{Data/daisy_ensemble.csv.gz}.

\paragraph{Table 4 of the article.} \emph{The 17 experiments of the study and their prepared designs in the workbench. Hours are the sum of the Daisy wall-clock run times of the study (\texttt{run\_seconds}) and the estimate for the compact design (runs $\times$ the median run time of the experiment). The experiments marked $-$ are not offered as designs; their published runs are bundled.}

\textbf{How it is made.} Table typed in the manuscript from the result files of the script. Script: \nolinkurl{p4_design_sizes.js}. Input: \nolinkurl{Data/daisy_ensemble.csv.gz}.

\paragraph{Table 5 of the article.} \emph{The Richards solver of the workbench against an independent reference (method of lines, scipy BDF, 0.5\,cm nodes). A 100\,cm column, initial suction 300\,cm, 200\,mm of rain on day 1, five dry days, free drainage. $D$ is the drainage below 1\,m and $S$ the final storage in mm, for the reference (ref) and the workbench (wb, 0.5\,cm nodes); the spread (mm) is the range over node spacings of 4, 2, 1 and 0.5\,cm.}

\textbf{How it is made.} Table typed from the result file (5 of its numbers found there). Script: \nolinkurl{Manuscripts/reproducibility/run_all.sh}. Input: \nolinkurl{Data/daisy_ensemble.csv.gz}.

\paragraph{Table 6 of the article.} \emph{Test questions of the planner and the factors that the offline reader returned (all of the expected factors, station and soil were present in every case).}

\textbf{How it is made.} Descriptive table (definitions, design, parameters or tests); no computation.

\paragraph{Table 7 of the article.} \emph{Mean seasonal nitrogen budget of the two bundled annual nitrogen logs as reported by the bridge (kg\,N\,ha$^{-1}$, ten seasons).}

\textbf{How it is made.} Table typed in the manuscript from the result files of the script. Script: \nolinkurl{(see the reproducibility appendix of the paper)}. Input: \nolinkurl{Data/daisy_ensemble.csv.gz}.

\paragraph{Table 8 of the article.} \emph{The three routes from Daisy to a decision.}

\textbf{How it is made.} Descriptive table (definitions, design, parameters or tests); no computation.

\paragraph{Table 9 of the article.} \emph{Deterministic phenotyping layer: maize in a warm and in a hot, dry season, seven genotypes, well-watered bench against a drought pulse (onset day 40, 28 days, 30\,\% of the plant-available water, rain-out shelter). Grain dry mass in model units; the penalty is the loss by the drought.}

\textbf{How it is made.} Table typed in the manuscript from the result files of the script. Script: \nolinkurl{(see the reproducibility appendix of the paper)}. Input: \nolinkurl{Data/daisy_ensemble.csv.gz}.

\paragraph{Table 10 of the article.} \emph{What this article adds, item by item: what was established before, what is new here, and the evidence in this article.}

\textbf{How it is made.} Descriptive table (definitions, design, parameters or tests); no computation.
